(1) 开路电势\(\mathcal E=\alpha\Delta T\)，串联欧姆定律给\(I=\alpha\Delta T/(R+R_L)=\alpha\Delta T/[R(1+m)]\)。输出\(P=I^2R_L=\alpha^2\Delta T^2 m/[R(1+m)^2]\)。对\(m>0\)求导，\(d[m/(1+m)^2]/dm=(1-m)/(1+m)^3\)，故\(m=1\)时功率最大，\(P_{\max}=\alpha^2\Delta T^2/(4R)\)。
(2) 热端吸热\(Q_h=\alpha T_h I-\tfrac12I^2R+K\Delta T\)，冷端放热\(Q_c=\alpha T_c I+\tfrac12I^2R+K\Delta T\)。两端Peltier热与Joule热符号相反，故\(Q_h-Q_c=\alpha\Delta T\,I-I^2R=I(\alpha\Delta T-IR)=I^2R_L=P\)，与电学输出严格相等，核对了本模型的能量守恒。
(3) 将\(I=\alpha\Delta T/[R(1+m)]\)代入\(\eta=P/Q_h\)并令\(Z=\alpha^2/(RK)\)、\(\bar T=(T_h+T_c)/2\)，先约去公共因子\(\alpha^2\Delta T/R\)后有\(\eta=\Delta T\,m/\{T_h(1+m)-\Delta T/2+(1+m)^2/Z\}\)。对\(m\)求导并令导数零：分母\(D=(1+m)^2/Z+T_h(1+m)-\Delta T/2\)，要求\(D-mD'=0\)，即\((1-m^2)/Z+T_h-\Delta T/2=0\)，所以\(m_{\eta,\max}=\sqrt{1+Z\bar T}\)。
(4) 将\(m_{\eta,\max}\)代回效率并整理，得\(\eta_{\max}=(\Delta T/T_h)[\sqrt{1+Z\bar T}-1]/[\sqrt{1+Z\bar T}+T_c/T_h]\)。\(Z\bar T\to0\)时\(\eta_{\max}\to0\)；\(Z\bar T\to\infty\)时逼近Carnot效率\(\eta_C=\Delta T/T_h\)，但有限\(Z\)下\(K>0\)及电阻焦耳耗散使\(\eta<\eta_C\)。最大效率负载通常\(m>1\)，与最大功率条件\(m=1\)不同。
